\chapter{Integrated discussion}
\label{chap:discussion}

The experiments separate three questions that are often conflated in method comparisons:
whether the designed order is recovered, whether different implementations solve the same
discrete problem, and whether the resulting accuracy is obtained at a useful cost.  The
combined evidence now permits a direct answer.  Higher order is valuable in this code, but
only when the reconstruction or predictor is paired with hardware that can amortise its
additional work.  Among the tested combinations, WENO-3 with HLLC on CUDA gives the most
consistent balance of accuracy, robustness, and runtime.

\section{When does higher order justify its cost?}

On smooth advection, WENO-3 and ADER-3 recover third order while WENO-2 recovers second
order.  This formal distinction remains visible in the practical \(600^2\) experiment:
WENO-3 and ADER-3 reduce \(L_1(\rho)\) by factors of 27.2 and 28.0 relative to WENO-2.
The two third-order errors are almost identical, so runtime decides between them.  On CUDA,
WENO-3 costs 2.71 s compared with 2.08 s for WENO-2, whereas ADER-3 costs 10.51 s.
WENO-3 therefore purchases most of the available smooth-case accuracy improvement for a
small additional GPU cost.  ADER-3 verifies that the one-step construction is accurate and
highly parallel, but its present predictor does not improve the error enough to recover its
larger absolute cost.

The discontinuous evidence changes the criterion.  Formal order is no longer achieved at
shocks, so contact resolution, positivity, and local oscillation become more important
than the smooth convergence slope.  WENO-3 again gives the strongest compromise: it is
less diffusive than WENO-2 around narrow structures, remains close to WENO-2 in the
far-field states, and costs substantially less than ADER-3.  The ADER-3 fallback prevents
catastrophic failure and preserves positive states, but the predictor/fallback transition
introduces the largest local backend sensitivity and does not give a uniformly cleaner
field.  Thus the main advantage of ADER-3 in the current code is its demonstration of a
robust, one-step third-order route rather than superiority on every test.

\section{What changes when HLL is replaced by HLLC?}

HLLC is the authoritative flux because the tested Euler problems contain material contacts
that are absent from the two-wave HLL model.  The matched comparison in
Figure~\ref{fig:hll-hllc-n600-tradeoff} also shows that this decision is supported by the
computed data rather than by theory alone.  At \(600^2\), changing from HLL to HLLC alters
the WENO-2 density error by less than 1\%, but lowers the WENO-3 and ADER-3 errors by about
66\%.  A more accurate reconstruction can therefore expose the benefit of the contact wave
more clearly than the second-order reference.

The runtime comparison is deliberately interpreted as implementation evidence.  HLLC is
not algebraically cheaper than HLL, yet the stored WENO HLLC paths are faster because the
complete kernels differ in branching, memory traffic, and compiler output.  ADER-3, for
which the predictor dominates the flux cost, changes by only 4.8\% on CUDA.  The defensible
conclusion is consequently that HLLC improves the measured accuracy and is not prohibitively
expensive in this implementation. It is not that HLLC must always run faster than HLL.

\section{Suitability of each backend}

Serial execution remains essential as the transparent reference implementation, but its
cost rises rapidly for the third-order methods.  OpenMP is most useful in an intermediate
regime.  It avoids CUDA's fixed launch and device-management costs on small grids and gives
useful speedup on a shared-memory CPU, but its efficiency eventually flattens.  The phase
profile explains why: once the stencil loops are parallel, workspace allocation, reductions,
updates, barriers, and memory bandwidth occupy a larger fraction of the run.  At 16 threads
and \(100^2\), the measured efficiencies are 26.3\%, 33.3\%, and 89.4\% for WENO-2,
WENO-3, and ADER-3.  The compute-heavy ADER predictor supplies enough work to use the
threads effectively, whereas the lighter WENO loops encounter overhead sooner.

CUDA is the appropriate backend for large grids and for the two third-order methods.  At
\(600^2\), smooth speedups reach 38.43 for WENO-2, 83.84 for WENO-3, and 392.01 for
ADER-3.  Phase profiling attributes this gain to the spatial operator: its CUDA speedup is
about 204 for both WENO methods and 1049 for ADER-3.  Acceleration therefore shifts the
bottleneck rather than removing it.  CFL reduction and state update become visible once the
dominant stencil or predictor work has been shortened.  Future optimization should target
those residual phases and reduce synchronization, rather than only tuning an already fast
flux kernel.

\begin{table}[H]
\centering
\small
\renewcommand{\arraystretch}{1.30}
\begin{tabular}{p{0.25\textwidth}p{0.26\textwidth}p{0.40\textwidth}}
\toprule
Objective & Recommended route & Reason \tabularnewline
\midrule
Low-cost reference or modest accuracy & WENO-2 on serial/OpenMP & Lowest method cost, useful before CUDA overhead is amortised. \tabularnewline
High smooth accuracy & WENO-3/HLLC on CUDA & Near-third-order error at a runtime close to CUDA WENO-2. \tabularnewline
Contact-rich discontinuous flow & WENO-3/HLLC on CUDA & Better local resolution than WENO-2 without ADER-3's present cost and sensor sensitivity. \tabularnewline
One-step high-order research route & ADER-3/HLLC on CUDA & Correct smooth order and very large parallel speedup, but further predictor/fallback optimization is required. \tabularnewline
\bottomrule
\end{tabular}
\caption{Recommended method/backend combinations within the tested problem range.}
\label{tab:recommendations}
\end{table}

\section{Limits of the evidence}

The conclusions are bounded by the experiment design.  The Gaussian problem supplies an
exact error measure, but the two-dimensional discontinuous cases do not. Their assessment
therefore combines positivity, backend agreement, field structure, and one-dimensional
slices.  Timings are best repeated values on a shared node and are intended to suppress
transient interference, not to reproduce an exclusively reserved machine.  The phase
profiler inserts CUDA synchronizations to assign kernel time and is diagnostic rather than
the source of the headline runtimes.

The official large-grid study ends at \(5000^2\).  It shows continued accuracy reduction
and increasing CUDA cost, but no single sharp saturation point that can be defended from
runtime slope alone.  Larger exploratory points were more sensitive to memory pressure and
shared-machine variability and were therefore not promoted into the canonical results.
Finally, the ADER-3 troubled-interface decision is intentionally nonlinear.  Small backend
roundoff can change a branch near its threshold, so a localized difference larger than
roundoff is possible even when both solutions remain conservative, positive, and physically
equivalent at the resolved scale.
